%% ma: L12-B08-0043
%% lop: 12
%% chuong: 2
%% bai: 8
%% dang: 12.8.29
%% loai: TN4
%% muc: NB
%% dap_an: "A"
%% nguon: "LaTeX giáo viên gửi 21/10/2026 (M05, Vectơ và hệ trục tọa độ trong không gian, 107 câu), Phần I Câu 36"
%% kien_thuc: "Bài 7 (tọa độ vectơ theo i, j, k), Bài 8 (biểu thức tọa độ phép cộng, nhân với số)"
%% trang_thai: chinh-thuc
%% ghi_chu: "Giáo viên đã duyệt 21/10/2026. Phần I Câu 48 của tài liệu (u - v với u, v cho theo i, j, k) cùng mô hình, chỉ đổi số nên đã bỏ"
\begin{ex}
Trong không gian $Oxyz$, cho vectơ $\vec{a}=2\vec{i}-3\vec{j}+\vec{k}$ và $\vec{b}=(1;2;-2)$. Tọa độ của vectơ $\vec{u}=\vec{a}+2\vec{b}$ là
\choice{\True $(4;1;-3)$}{$(3;1;-1)$}{$(4;-1;3)$}{$(0;-7;5)$}
\loigiai{$\vec{a}=(2;-3;1)$, $2\vec{b}=(2;4;-4)$ nên $\vec{u}=(4;1;-3)$.}
\end{ex}
